\subsection{DIRECT SUMS OF SUBMODULES}

DEFINITION 6. An A-module E is said to be the direct sum of a family \((\mathbf{M}_{\iota})_{\iota \in I}\) of submodules of E if the canonical mapping \(\bigoplus_{\iota \in I} \mathbf{M}_{\iota} \to \mathbf{E}\) (no. 6) is an isomorphism.

It amounts to the same to say that every \(x \in \mathbf{E}\) can be written in a unique way in the form \(x = \sum_{i \in I} x_i\) , where \(x_i \in \mathbf{E}_i\) for all \(i \in I\) ; the element \(x_i\) thus corresponding to \(x\) is called the component of \(x\) in \(\mathbf{E}_i\) ; the mapping \(x \mapsto x_i\) is linear.

Remark. (1) Let \((\mathbf{M}_{\iota})_{\iota \in \mathbf{I}}, (\mathbf{N}_{\iota})_{\iota \in \mathbf{I}}\) be two families of submodules of a module \(\mathbf{E}\) , with the same indexing set; suppose that \(\mathbf{E}\) is both the direct sum of the family \((\mathbf{M}_{\iota})\) and of the family \((\mathbf{N}_{\iota})\) and that \(\mathbf{N}_{\iota} \subset \mathbf{M}_{\iota}\) for all \(\iota \in \mathbf{I}\) . Then \(\mathbf{N}_{\iota} = \mathbf{M}_{\iota}\) for all \(\iota \in \mathbf{I}\) , as follows immediately from no. 6, Corollary 1 to Proposition 7 applied to the canonical injections \(f_{\iota}: \mathbf{N}_{\iota} \to \mathbf{M}_{\iota}\) .

PROPOSITION 11. Let \((\mathbf{M}_{i})_{i\in I}\) be a family of submodules of an A-module E. The following properties are equivalent:

\begin{enumerate}
\def\labelenumi{(\alph{enumi})}
\item
  The submodule \(\sum_{i\in I}\mathbf{M}_i\) is the direct sum of the family \((\mathbf{M}_i)_{i\in I}\) .
\item
  The relation \(\sum_{i\in I}x_i = 0\) where \(x_{i}\in \mathbf{M}_{i}\) for all \(i\in I\) , implies \(x_{i} = 0\) for all \(i\in I\)
\item
  For all \(\kappa \in \mathbf{I}\) , the intersection of \(\mathbf{M}_{\kappa}\) and \(\sum_{i \neq \kappa} \mathbf{M}_i\) reduces to 0.
\end{enumerate}

It is immediate that (a) and (b) are equivalent, since the relation \(\sum_{i} x_{i} = \sum_{i} y_{i}\) is equivalent to \(\sum_{i} (x_{i} - y_{i}) = 0\) . On the other hand, by virtue of Definition 6, (a) implies (c) by the uniqueness of the expression of an element of \(\bigoplus_{i \in I} M_{i}\) as a direct sum of elements \(x_{i} \in M_{i}\) . Finally, the relation \(\sum_{i} x_{i} = 0\) , where \(x_{i} \in M_{i}\) for all \(i\) , can be written, for all \(\kappa \in I\) , \(x_{\kappa} = \sum_{i \neq \kappa} (-x_{i})\) ; condition (c) then implies \(x_{\kappa} = 0\) for all \(\kappa \in I\) , hence (c) implies (b).

DEFINITION 7. An endomorphism \(e\) of an A-module \(\mathbf{E}\) is called a projector if \(e \circ e = e\) (in other words, if \(e\) is an idempotent in the ring \(\operatorname{End}(\mathbf{E})\) ). In \(\operatorname{End}(\mathbf{E})\) a family \((e_{\lambda})_{\lambda \in L}\) of projectors is called orthogonal if \(e_{\lambda} \circ e_{\mu} = 0\) for \(\lambda \neq \mu\) .

PROPOSITION 12. Let E be an A-module.

\begin{enumerate}
\def\labelenumi{(\roman{enumi})}
\item
  If \(\mathbf{E}\) is the direct sum of a family \((\mathbf{M}_{\lambda})_{\lambda \in \mathbf{L}}\) of submodules and, for all \(x \in \mathbf{E}\) , \(e_{\lambda}(x)\) is the component of \(x\) in \(\mathbf{M}_{\lambda}\) , \((e_{\lambda})\) is an orthogonal family of projectors such that \(x = \sum_{\lambda \in \mathbf{L}} e_{\lambda}(x)\) for all \(x \in \mathbf{E}\) .
\item
  Conversely, if \((e_{\lambda})_{\lambda \in \mathbf{L}}\) is an orthogonal family of projectors in \(\operatorname{End}(\mathbf{E})\) such that \(x = \sum_{\lambda \in \mathbf{L}} e_{\lambda}(x)\) for all \(x \in \mathbf{E}\) , \(\mathbf{E}\) is the direct sum of the family of submodules \(\mathbf{M}_{\lambda} = e_{\lambda}(\mathbf{E})\) .
\end{enumerate}

Property (i) follows from the definitions and (ii) is a special case of no. 6, Corollary 2 to Proposition 6, applied to the canonical injections \(\mathbf{M}_{\lambda} \to \mathbf{E}\) and the mappings \(e_{\lambda}: \mathbf{E} \to \mathbf{M}_{\lambda}\) .

Note that when \(\mathbf{L}\) is finite the condition \(x = \sum_{\lambda \in \mathbf{L}} e_{\lambda}(x)\) for all \(x \in \mathbf{E}\) may also be written in \(\operatorname{End}(\mathbf{E})\) .

\[
1 _ {\mathbf {E}} = \sum_ {\lambda \in L} e _ {\lambda}.\tag{31}
\]

COROLLARY. For every projector \(e\) of \(\mathbf{E}\) , \(\mathbf{E}\) is the direct sum of the image \(\mathbf{M} = e(\mathbf{E})\) and the kernel \(\mathbf{N} = e^{-1}(0)\) of \(e\) ; for all \(x = x_1 + x_2 \in \mathbf{E}\) with \(x_1 \in \mathbf{M}\) and \(x_2 \in \mathbf{N}\) , \(x_1 = e(x)\) ; 1 - \(e\) is a projector of \(\mathbf{E}\) of image \(\mathbf{N}\) and kernel \(\mathbf{M}\) .

\((1 - e)^2 = 1 - 2e + e^2 = 1 - e\) in \(\operatorname{End}(\mathbf{E})\) and hence \(1 - e\) is a projector; as also \(e(1 - e) = (1 - e)e = e - e^2 = 0\) , \(\mathbf{E}\) is the direct sum of the images \(\mathbf{M}\) and \(\mathbf{N}\) of \(e\) and \(1 - e\) by Proposition 12. Finally, for all \(x \in \mathbf{E}\) , the relation \(x \in \mathbf{M}\) is equivalent to \(x = e(x)\) ; for \(x = e(x)\) implies by definition \(x \in \mathbf{M}\) and, conversely, if \(x = e(x')\) with \(x' \in \mathbf{E}\) , then \(e(x) = e^2(x') = e(x') = x\) ;

this shows therefore that M is the kernel of 1 - e and, exchanging the roles of e and 1 - e, it is similarly seen that N is the kernel of e.

Remark. (2) Let E, F be two A-modules such that E is the direct sum of a finite family \((\mathbf{M}_{i})_{1\leqslant i\leqslant m}\) of submodules and F the direct sum of a finite family \((\mathbf{N}_{j})_{1\leqslant j\leqslant n}\) of submodules. Then it is known (no. 6, Corollary 1 to Proposition 6) that \(\operatorname{Hom}_{\mathbf{A}}(\mathbf{E},\mathbf{F})\) is canonically identified with the product \(\prod_{i,j}\operatorname{Hom}_{\mathbf{A}}(\mathbf{M}_{i},\mathbf{N}_{j})\) ; to be precise, to a family \((u_{ji})\) , where \(u_{ji}\in\operatorname{Hom}_{\mathbf{A}}(\mathbf{M}_{i},\mathbf{N}_{j})\) there corresponds the linear mapping \(u:E\to F\) defined as follows. It suffices to define the restriction of u to each of the \(M_{i}\) and for each \(x_{i}\in M_{i}\) ,

\[
u (x _ {i}) = \sum_ {j = 1} ^ {n} u _ {j i} (x _ {i}).
\]

Now let \(\mathbf{G}\) be a third A-module, the direct sum of a finite family \((\mathbf{P}_k)_{1 \leqslant k \leqslant p}\) of submodules; let \(v\) be a linear mapping of \(\mathbf{F}\) into \(\mathbf{G}\) and let \((v_{kj}) \in \prod_{j,k} \operatorname{Hom}_{\mathbf{A}}(\mathbf{N}_j, \mathbf{P}_k)\) be the family corresponding to it canonically. For all \(x_i \in \mathbf{M}_i\) ,

\[
v (u (x _ {i})) = \sum_ {j = 1} ^ {n} v (u _ {j i} (x _ {i})) = \sum_ {k = 1} ^ {p} \sum_ {j = 1} ^ {n} v _ {k j} (u _ {j i} (x _ {i})).
\]

Thus it is seen that if we write

\[
w _ {k i} = \sum_ {j = 1} ^ {n} v _ {k j} \circ u _ {j i} \in \operatorname{Hom} _ {\mathbf {A}} (\mathrm{M} _ {i}, \mathrm{P} _ {k})\tag{32}
\]

the family \((w_{ki})\) corresponds canonically to the composite linear mapping \(w = v \circ u\) from \(\mathbf{E}\) to \(\mathbf{G}\) (cf.~§ 10, no. 5).

